% Compile twice: lualatex 20001315.tex
\documentclass[11pt]{article}
\usepackage[margin=25mm]{geometry}
\usepackage{fontspec,unicode-math}
\setmainfont{Latin Modern Roman}
\setmathfont{Latin Modern Math}
\usepackage{amsmath,amssymb,mathtools}
\usepackage{xurl,hyperref}
\hypersetup{colorlinks=true,urlcolor=blue,linkcolor=blue}
\setcounter{secnumdepth}{0}
\setlength{\parindent}{0pt}
\setlength{\parskip}{5pt}
\setlength{\emergencystretch}{3em}
\newcommand{\sref}[2]{\hyperlink{source:#1:#2}{[S#2]}}
\newcommand{\cataloguescope}[1]{\paragraph{Source question.}#1}
\newcommand{\floor}[1]{\left\lfloor#1\right\rfloor}
\begin{document}
% UnsolvedMath ID 20001315; source catalogue code AIM-CRYPTOGRAPHY-0002
\setcounter{section}{20001314}
\section{Regular-orbit equivalence and affine orbit-size certificates for SDLP}\label{20001315}
{\small\sffamily UnsolvedMath ID 20001315 · Computer Science}
\subsection{Definitions and mathematical statement}

\paragraph{Shared notation.}Standard number systems are denoted by $\mathbb N,\mathbb Z,\mathbb Q,\mathbb R,\mathbb C$; index ranges, computational models and specialized notation follow the cited problem source and the definitions below.


A semidirect product combines a group and an action by automorphisms. In a semidirect discrete-logarithm problem, a projected component of $(g,\phi)^n$ is given and the exponent n is sought; the corresponding Diffie--Hellman-type task combines projected secret powers. A composition series is a finite normal series with simple factors.
\paragraph{Statement recorded in the supplied source.}
Cryptanalysis of SDLP

Various problems relating to the cryptanalysis, both quantum and classical, of SDLP remain open.

On the quantum side the question of quantum equivalence of SDLP and the CDH-type problem underpinning semidirect product key exchange (SCDH) remains unproven, though we have strong cause to suspect their equivalence given the quantum equivalence of the heavily related vectorisation and parallelisation problems. The decomposition attack of Imran and Ivanyos, providing in some cases a reduction of SDLP to the abelian hidden subgroup problem, relies on the quantum ability to quickly compute a composition series of a solvable group. Whether similar quantum speedups in the computation series exist for non-solvable groups remains open.

On the classical side, we would like to develop a kind of Sylow theory for the projected sets implied by the semidirect product structure. It would also be useful to develop a classical method of computing the size of the acting group when one considers the semigroup variant of SDLP (currently we have only quantum methods). \sref{20001315}{2}
\paragraph{Scope qualification.}
The source bundles several questions and includes a semigroup variant; the actual group encodings and action are not fixed, and no one uniform equivalence is assumed.
\paragraph{Source status.}
The supplied catalogue labels this record partially solved. Its source-status metadata was checked against the supplied snapshot on 2026-10-11. This import does not independently establish current open status or mathematical correctness.
\subsection{Short English statement}
What quantum and classical reductions, composition-series algorithms, and orbit-size methods are available for semidirect discrete-logarithm problems?
\subsection{Sources}
\begin{itemize}
\item[\hypertarget{source:20001315:1}{S1}] Ulam AI and the UnsolvedMath Contributors, UnsolvedMath: A Curated Collection of Open Mathematics Problems, record 20001315 (AIM-CRYPTOGRAPHY-0002). Catalogue curation and metadata: CC BY 4.0. The supplied source question is retained; definitions and the English summary are editorial additions.
\url{https://huggingface.co/datasets/ulamai/UnsolvedMath}.
\item[\hypertarget{source:20001315:2}{S2}] AIM Workshop Problem Lists, original problem source linked by catalogue record AIM-CRYPTOGRAPHY-0002
\url{http://aimpl.org/postquantgroup/1/}.
\end{itemize}
\paragraph{Preparation and review.}
Original question attributed to the cited source and UnsolvedMath contributors. Codex (GPT-6; exact serving revision not exposed) assisted with definitions, English summaries, source inspection and formatting. No new human mathematical authorship or independent proof verification is asserted. Canonical ID 20001315 and source code AIM-CRYPTOGRAPHY-0002 checked against the complete supplied registry and existing repository statements; known flagged repetitions and obvious identical conjectures excluded. Full mathematical-equivalence/literature review has not been performed. Statement preparation and final review followed contribution prompts C and D; automated checks verify structure and identity, not open status.
\label{end:20001315}
\end{document}
